\documentclass{article}
\usepackage{graphicx} % Required for inserting images
\usepackage{parskip}
\usepackage{url}
\usepackage{longtable}

\title{Tesi - Alessandro Porcheddu}
\author{Alessandro Ricci}
\date{June 2026}

\begin{document}

\maketitle

% \begin{abstract}
% Questa tesi presenta il progetto di una centralina elettronica programmabile
% per motori monocilindrici a carburatore. Il lavoro considera la centralina non
% come un dispositivo isolato, ma come il nodo real-time di un sistema che
% comprende acquisizione sensoristica, controllo motore, telemetria, interfaccia
% utente ed elaborazione dei dati. La ECU acquisisce e qualifica i segnali del
% motore, calcola le strategie di controllo, temporizza accensione e attuatori e
% rende osservabili misure, decisioni ed eventi attraverso una telemetria locale.
% Il sistema comprende inoltre la gestione delle mappe, la registrazione delle
% sessioni, l'aggiornamento del firmware e un digital twin lato server.
% \end{abstract}

% \tableofcontents

\section{Introduzione e contesto applicativo}
\label{sec:introduzione}

\subsection{Dominio applicativo}

Le centraline elettroniche, comunemente indicate con il termine Electronic Control Unit (ECU), rappresentano uno degli elementi fondamentali dei moderni sistemi automotive. Una ECU è un sistema embedded dedicato al controllo di specifiche funzioni del veicolo e, nel caso del motore, ha il compito di acquisire informazioni dai sensori, elaborarle secondo opportune strategie di controllo e generare i comandi necessari agli attuatori.

Nel corso degli anni, le ECU si sono evolute da semplici dispositivi di controllo a piattaforme software sempre più sofisticate. Nelle applicazioni automotive moderne e nelle competizioni motociclistiche di alto livello, le centraline non si limitano a gestire l'accensione o l'iniezione, ma integrano funzionalità avanzate quali telemetria, diagnostica, registrazione dei dati, configurazione dei parametri operativi e sistemi di protezione del motore. Questi strumenti consentono di osservare il comportamento del sistema, adattarne il funzionamento alle esigenze dell'applicazione e intervenire tempestivamente in presenza di condizioni anomale.

Tale evoluzione non ha però interessato in modo uniforme tutti i settori. Nel contesto dei motori monocilindrici a due tempi, molte delle soluzioni attualmente disponibili si basano ancora su tecnologie e architetture sviluppate diversi decenni fa. Le centraline impiegate in questo ambito offrono generalmente funzionalità limitate: la telemetria è spesso assente, gli strumenti diagnostici sono ridotti o inesistenti e le possibilità di configurazione risultano molto ristrette. In molti casi non è possibile modificare facilmente la curva di anticipo, adattare la risposta all'acceleratore o personalizzare il comportamento della centralina in funzione delle caratteristiche del motore e delle esigenze dell'utilizzatore.

Anche gli strumenti dedicati alla sicurezza e alla salvaguardia del propulsore risultano spesso limitati. L'assenza di meccanismi di monitoraggio e protezione rende più difficile individuare condizioni operative potenzialmente dannose e intervenire prima che possano causare guasti o degrado delle prestazioni. Di conseguenza, molte attività di calibrazione e verifica continuano a dipendere da prove empiriche e dall'esperienza dell'operatore.

\subsection{Idea di progetto}

A partire dal contesto descritto, l'idea alla base del progetto è realizzare una ECU custom per motori monocilindrici a due tempi che non si limiti alla sola generazione dei segnali di controllo, ma che integri anche funzionalità normalmente associate a sistemi più evoluti. L'obiettivo è trasferire in questo ambito alcuni concetti ormai consolidati in altri settori automotive e motorsport, come la telemetria, la diagnostica, la configurazione dei parametri operativi e la protezione del motore.

La centralina viene quindi considerata come una piattaforma embedded completa, capace di acquisire informazioni dal motore, elaborarle in tempo reale e renderle disponibili per l'osservazione e la validazione del comportamento del sistema. Questa impostazione permette di superare il modello tradizionale della ECU come dispositivo chiuso e difficilmente ispezionabile, introducendo invece un'architettura più trasparente, configurabile ed estendibile.

Dal punto di vista funzionale, il progetto prevede che la ECU sia in grado di gestire i segnali provenienti dai sensori, determinare lo stato operativo del motore e generare i comandi necessari agli attuatori rispettando i vincoli temporali imposti dal funzionamento del propulsore. Parallelamente, il sistema deve poter raccogliere e trasmettere dati di funzionamento, rilevare condizioni anomale, conservare configurazioni persistenti e fornire strumenti utili alla verifica delle strategie implementate.

Un aspetto centrale dell'idea progettuale è la separazione tra controllo del motore e funzionalità di supporto. Il controllo real-time deve mantenere la priorità assoluta, mentre telemetria, diagnostica, configurazione e registrazione dei dati devono essere progettate in modo da non compromettere l'affidabilità della centralina. Questa distinzione consente di introdurre funzionalità avanzate senza trasformarle in un rischio per il corretto funzionamento del motore.

La ECU proposta non viene quindi intesa esclusivamente come un componente elettronico dedicato all'accensione o alla gestione di singoli attuatori, ma come il nucleo di un sistema più ampio. Attraverso l'integrazione con strumenti software esterni, essa può rendere il motore osservabile durante il funzionamento, permettere la modifica controllata dei parametri e supportare attività di analisi, calibrazione e validazione. In questo modo il progetto mira a costruire una base tecnica su cui sviluppare progressivamente funzionalità più avanzate, mantenendo una struttura coerente e adatta a un'applicazione embedded real-time.

\subsection{Riferimenti architetturali e stato dell'arte}

L'indipendenza tra logica applicativa e dettagli della piattaforma è un tema
consolidato nel software automotive. AUTOSAR Classic, ad esempio, separa
applicazione, runtime e software di base; l'applicazione è in larga parte
indipendente dall'hardware e accede ai servizi mediante interfacce
definite~\cite{autosar_classic}. L'architettura adotta una separazione calibrata
sulle dimensioni del sistema: il codice che interpreta un sensore non conosce
il driver ESP-IDF che ha prodotto il campione e i servizi applicativi dipendono
da contratti stabili.

Un secondo riferimento è la validazione hardware-in-the-loop (HIL). In questo
approccio il software esegue sul controllore reale mentre una sorgente esterna
riproduce il comportamento dell'ambiente fisico. Abboush, Knieke e Rausch
mostrano come HIL e iniezione di guasti possano sostenere la validazione real-time
del software automotive~\cite{abboush_hil}. Il simulatore del sistema genera il
pick-up e gli ingressi analogici o digitali e li presenta alla ECU attraverso
le normali interfacce elettriche. Le prove sono quindi ripetibili senza
modificare la logica sottoposta a verifica.

Il terzo riferimento riguarda il digital twin. La letteratura sui motori a
combustione interna associa il digital twin a monitoraggio, diagnosi e modelli
predittivi alimentati dai dati del sistema fisico~\cite{tran_dt}. Lavori più
recenti sui motori diesel combinano modelli fisici e tecniche di apprendimento
per la stima dei parametri e il monitoraggio dello stato~\cite{nath_dt}. Nel
sistema sviluppato il twin mantiene la rappresentazione persistente della ECU,
delle revisioni installate e delle sessioni registrate. Su questa base esegue
replay, confronto tra mappe e analisi diagnostica, mantenendo separati dati
misurati, risultati ricalcolati e valori simulati.

\subsection{Obiettivo e contributo della tesi}

Il contributo è duplice. Sul piano progettuale viene definito un sistema
completo, dai segnali fisici alla conservazione delle sessioni, con responsabilità
separate. Sul piano sperimentale il sistema integra acquisizione e modellazione
degli ingressi, controllo dell'accensione e degli attuatori, gestione delle
mappe, pubblicazione di snapshot ed eventi, trasporto WebSocket, interfaccia web
locale e servizi remoti per sessioni, aggiornamenti e digital twin.

La validazione segue l'intera catena: verifica che gli ingressi vengano
acquisiti e qualificati, che le decisioni siano applicate con la temporizzazione
richiesta, che la telemetria non blocchi il controllo e che sessioni e
configurazioni rimangano tracciabili.

\section{Analisi dei requisiti}
\label{sec:requisiti}

\subsection{Attori e confine del sistema}

Il sistema principale è costituito dalla ECU, responsabile dell'acquisizione dei segnali, dell'elaborazione delle strategie di controllo e dell'attuazione dei comandi verso il motore.

Gli attori esterni principali sono quattro: il motore, il pilota, il tecnico e il servizio remoto. Il motore rappresenta il processo fisico da controllare: da esso vengono acquisite le grandezze necessarie alla ricostruzione dello suo stato, come il regime di rotazione, la posizione angolare dell'albero motore, la posizione della farfalla, la temperatura dei gas di scarico e gli eventuali segnali associati alla detonazione. Allo stesso tempo, il motore è anche il destinatario finale delle azioni calcolate dalla ECU, come il comando di accensione della scintilla, la gestione del power jet e il controllo della valvola di scarico.

Il pilota interagisce con il sistema attraverso comandi diretti, legati alla guida del veicolo. Alcuni di questi segnali hanno carattere immediato e devono essere gestiti con priorità elevata, come la richiesta proveniente dal quick shifter o le variazioni rapide dell'apertura farfalla. Altri comandi, come la selezione della mappa o l'attivazione di strategie specifiche, non richiedono necessariamente la stessa latenza, ma devono comunque essere acquisiti e applicati in modo coerente con lo stato corrente della centralina.

Il tecnico rappresenta invece l'utente incaricato della configurazione, della calibrazione e della diagnosi del sistema. La sua interazione avviene tramite computer o smartphone collegandosi con un browser alla webapp fornita dalla ECU. Attraverso tale interfaccia il tecnico può osservare la telemetria, modificare parametri e mappe, verificare lo stato degli ingressi e degli attuatori, ed effettuare attività di test o manutenzione.

Il servizio remoto ha un ruolo distinto rispetto al ciclo di controllo. Esso può essere utilizzato per conservare dati, revisioni delle configurazioni, informazioni diagnostiche, statistiche sulla telemetria ed eventualmente per distribuire aggiornamenti software. La disponibilità della rete o del server non costituisce una condizione necessaria per il funzionamento della ECU.

\subsection{Requisiti funzionali}

\begin{description}
    \item[RF-001.] Il sistema deve acquisire un riferimento di posizione e
          velocità del motore per consentire il controllo sincronizzato delle
          funzioni motore.

    \item[RF-002.] Il sistema deve determinare lo stato dinamico del motore,
          inclusi regime di rotazione, variazioni di velocità e stato di
          sincronizzazione.

    \item[RF-003.] Il sistema deve validare la coerenza dei segnali utilizzati
          per il controllo motore e rilevare condizioni non plausibili.

    \item[RF-004.] Il sistema deve consentire il funzionamento delle strategie
          motore solo quando le informazioni necessarie risultano affidabili.

    \item[RF-005.] Il sistema deve calcolare il comando di accensione sulla base
          dello stato motore e delle strategie configurate.

    \item[RF-006.] Il sistema deve supportare la configurazione dei parametri
          necessari alla sincronizzazione e al controllo del motore.

    \item[RF-007.] Il sistema deve acquisire i principali parametri operativi
          del motore attraverso sensori dedicati.

    \item[RF-008.] Il sistema deve normalizzare, validare e rendere disponibili
          i dati provenienti dai sensori alle funzioni di controllo.

    \item[RF-009.] Il sistema deve supportare procedure di calibrazione dei
          sensori configurabili dall'utente.

    \item[RF-010.] Il sistema deve applicare strategie di fallback in caso di
          indisponibilità o invalidità dei dati sensoriali.

    \item[RF-011.] Il sistema deve gestire configurazioni persistenti e
          ripristinarle automaticamente all'avvio.

    \item[RF-012.] Il sistema deve validare configurazioni e strategie prima
          della loro applicazione operativa.

    \item[RF-013.] Il sistema deve utilizzare configurazioni di sicurezza
          predefinite in caso di dati mancanti, corrotti o non validi.

    \item[RF-014.] Il sistema deve consentire la configurazione e il tuning
          delle strategie motore tramite interfacce dedicate.

    \item[RF-015.] Il sistema deve fornire strumenti di monitoraggio in tempo
          reale dello stato del motore e della ECU.

    \item[RF-016.] Il sistema deve rendere disponibili informazioni operative e
          diagnostiche tramite telemetria.

    \item[RF-017.] Il sistema deve monitorare continuamente lo stato di salute
          dei sensori e delle funzioni interne.

    \item[RF-018.] Il sistema deve rilevare, classificare e registrare
          condizioni di errore e anomalie operative.

    \item[RF-019.] Il sistema deve acquisire parametri termici rilevanti per il
          monitoraggio dello stato del motore.

    \item[RF-020.] Il sistema deve fornire informazioni sullo stato termico del
          motore e sulle relative tendenze evolutive.

    \item[RF-021.] Il sistema deve rilevare fenomeni potenzialmente dannosi per
          il motore e renderli disponibili alle strategie di controllo.

    \item[RF-022.] Il sistema deve supportare ingressi dedicati all'interazione
          del pilota con le strategie motore.

    \item[RF-023.] Il sistema deve consentire la selezione e il cambio delle
          strategie operative durante il funzionamento.

    \item[RF-024.] Il sistema deve supportare la gestione di più mappe o profili
          di funzionamento selezionabili.

    \item[RF-025.] Il sistema deve fornire informazioni sullo stato corrente
          delle strategie attive e delle modalità operative.

    \item[RF-026.] Il sistema deve supportare funzionalità di diagnostica e
          manutenzione del sistema.

    \item[RF-027.] Il sistema deve consentire la simulazione delle principali
          condizioni operative per attività di verifica e validazione.

    \item[RF-028.] Il sistema deve supportare la raccolta e la consultazione dei
          dati necessari all'analisi delle prestazioni e dei fault.

    \item[RF-029.] Il sistema deve garantire la disponibilità delle informazioni
          necessarie al controllo, alla diagnostica e alla supervisione del
          motore.
\end{description}

\subsection{Requisiti non funzionali}

\begin{description}
    \item[RNF-001.] Le funzioni critiche di controllo motore devono operare con
          comportamento deterministico e prevedibile.

    \item[RNF-002.] Le funzioni critiche devono essere isolate dalle
          funzionalità non essenziali quali interfacce utente, comunicazioni e
          servizi ausiliari.

    \item[RNF-003.] Il sistema deve garantire la precisione temporale necessaria
          al corretto controllo del motore.

    \item[RNF-004.] Il sistema deve minimizzare la latenza tra acquisizione
          degli eventi critici e relativa elaborazione.

    \item[RNF-005.] Le attività non critiche non devono compromettere le
          prestazioni delle funzioni di controllo motore.

    \item[RNF-006.] I meccanismi di comunicazione e telemetria devono essere non
          bloccanti rispetto alle funzioni critiche.

    \item[RNF-007.] Il sistema deve avviarsi e operare in condizioni sicure fino
          alla verifica della propria integrità operativa.

    \item[RNF-008.] Il sistema deve essere in grado di rilevare e recuperare da
          condizioni di blocco o malfunzionamento software.

    \item[RNF-009.] Le configurazioni persistenti devono essere protette da
          corruzione e perdita di dati.

    \item[RNF-010.] Le strategie di controllo devono mantenere comportamento
          coerente indipendentemente dal carico delle altre funzioni del
          sistema.

    \item[RNF-011.] L'architettura software deve essere modulare e separare
          chiaramente responsabilità e domini funzionali.

    \item[RNF-012.] Il sistema deve essere progettato per favorire
          manutenibilità, evoluzione e riutilizzo dei componenti software.

    \item[RNF-013.] Le logiche di controllo devono essere verificabili e
          testabili indipendentemente dall'hardware finale.

    \item[RNF-014.] Il sistema deve supportare attività di simulazione e
          validazione ripetibili e automatizzabili.

    \item[RNF-015.] Le interfacce utente devono fornire una rappresentazione
          chiara dello stato operativo del sistema.

    \item[RNF-016.] Ogni informazione resa disponibile dal sistema deve essere
          accompagnata da indicatori di validità e affidabilità quando
          applicabile.

    \item[RNF-017.] Il sistema deve essere robusto rispetto a rumore, disturbi e
          condizioni operative non ideali.

    \item[RNF-018.] La telemetria deve garantire coerenza, tracciabilità e
          contestualizzazione dei dati trasmessi.

    \item[RNF-019.] In condizioni di risorse limitate, il sistema deve
          preservare prioritariamente le informazioni critiche per sicurezza e
          diagnostica.

    \item[RNF-020.] Gli eventi e i dati registrati devono mantenere il corretto
          ordine temporale.

    \item[RNF-021.] Le funzionalità diagnostiche avanzate devono essere
          configurabili e non interferire con il funzionamento operativo.

    \item[RNF-022.] L'accesso alle funzioni critiche deve essere protetto da
          adeguati meccanismi di autorizzazione e conferma.

    \item[RNF-023.] Il sistema deve adottare strategie di degradazione
          controllata in presenza di fault o perdita di informazioni.

    \item[RNF-024.] L'architettura deve essere estendibile e consentire
          l'introduzione di nuove funzionalità senza modifiche sostanziali ai
          componenti esistenti.
\end{description}

\subsection{Confine funzionale del sistema}

Il sistema copre acquisizione, stima dello stato, controllo dell'accensione,
attuazione ausiliaria, sicurezza, gestione delle mappe, telemetria, comandi,
registrazione delle sessioni, digital twin e aggiornamento OTA. Le funzioni
real-time risiedono sulla ECU; visualizzazione, conservazione e analisi sono
distribuite tra client e server.

La piattaforma di prova comprende il condizionamento dei segnali e il simulatore
necessari alla validazione HIL. La scheda dedicata, la qualificazione ambientale
e l'integrazione meccanica sul veicolo appartengono al processo di
industrializzazione e non modificano il confine funzionale dell'architettura.

\section{Progettazione del sistema}
\label{sec:progettazione}

\subsection{Macro-architettura}

L'architettura è composta da cinque sottosistemi. I primi quattro appartengono
al sistema operativo; il quinto esiste per verificarlo.

\begin{description}
    \item[Impianto fisico.] Comprende motore, sensori, condizionamento dei
          segnali e attuatori. Produce grandezze continue ed eventi angolari e
          riceve comandi elettrici.

    \item[Nodo ECU.] Acquisisce gli ingressi, stima lo stato del motore, applica
          controllo e sicurezza, comanda gli attuatori e pubblica telemetria. È
          l'unico sottosistema che deve rispettare le scadenze del motore.

    \item[Client locale.] Visualizza stato ed eventi, invia richieste validate
          e conserva temporaneamente i dati necessari all'esperienza utente.
          Comunica con la ECU sulla rete locale.

    \item[Server e digital twin.] Conserva identità, revisioni e sessioni. Offre
          confronto, replay e analisi storica. Non interviene nel controllo
          real-time.

    \item[Banco di simulazione.] Riproduce segnali fisici controllabili e
          ripetibili. Dal punto di vista della ECU sostituisce uno o più sensori,
          non il software che li interpreta.
\end{description}

Il client svolge anche il ruolo di gateway verso il server. Questa scelta evita
che la disponibilità di Internet diventi parte del ciclo di vita minimo della
ECU e consente di mantenere disattivata la comunicazione remota durante il
funzionamento critico. Le sessioni vengono accodate localmente e trasferite
quando il client e il collegamento remoto sono disponibili.

\subsection{Relazioni tra i sottosistemi}

Il simulatore e i sensori reali comunicano con la ECU mediante segnali elettrici:
livelli analogici, ingressi digitali, bus per convertitori e cattura degli eventi
di pick-up. Questo collegamento appartiene al dominio fisico e attraversa gli
stadi di protezione e condizionamento posti tra cablaggio e controllore.

ECU e client scambiano dati su un canale WebSocket locale. La direzione ECU--client
trasporta snapshot, eventi, contatori e capacità dichiarate dal firmware; la
direzione opposta trasporta richieste. Il protocollo bidirezionale non implica
responsabilità simmetriche: il client chiede una modifica, mentre la ECU decide
se sia valida e quando applicarla.

Il client inoltra al server sessioni e metadati attraverso una connessione non
real-time. Il server restituisce storico, artefatti approvati e risultati di
analisi. Per ogni sessione devono rimanere associati almeno firmware,
configurazione, mappa e calibrazione. Senza questi riferimenti due curve simili
potrebbero derivare da sistemi diversi e non sarebbero confrontabili.

\subsection{Architettura interna della ECU}

All'interno della ECU le dipendenze sono dirette dalla piattaforma verso il
dominio, non nella direzione opposta. La decomposizione scelta comprende:

\begin{itemize}
    \item \textbf{adattatori hardware}, che usano ADC, GPIO, SPI e periferiche di
          cattura e traducono i risultati in tipi indipendenti da ESP-IDF;
    \item \textbf{porte di acquisizione}, che definiscono le operazioni richieste
          dal dominio alle sorgenti fisiche o simulate;
    \item \textbf{oggetti e servizi di dominio}, che validano i campioni,
          calcolano stato di salute, debounce, timeout e grandezze derivate;
    \item \textbf{deposito sensoriale}, che pubblica istantanee coerenti
          (\emph{snapshot}) e code
          limitate di eventi;
    \item \textbf{consumatori}, costituiti da controllo motore, sicurezza,
          telemetria, diagnostica e registrazione;
    \item \textbf{runtime}, che crea i task, assegna priorità e collega le
          implementazioni alle porte.
\end{itemize}

L'isolamento ha una conseguenza pratica: il sensore TPS non decide un limite di
coppia e il quick shifter non porta direttamente a zero l'uscita CDI. Il primo
pubblica una lettura o un fallback esplicito; il secondo pubblica una richiesta
validata. La strategia di controllo e sicurezza conserva l'autorità sull'azione
finale. In caso contrario, una modifica al debounce di un ingresso digitale
potrebbe cambiare direttamente il comando dell'accensione.

\subsection{Proprietà dei dati e concorrenza}

Gli ingressi non hanno tutti la stessa semantica. Per TPS e temperature interessa
il valore valido più recente; un aggiornamento non trasmesso può essere
sostituito da uno più nuovo. Per cambio mappa, transizione di guasto e richiesta
quick shifter conta invece l'evento: sostituirlo silenziosamente cambierebbe la
storia osservata dal sistema. Il progetto usa quindi snapshot per lo stato e
code limitate per gli eventi.

Ogni produttore scrive solo i dati che possiede. I consumatori leggono copie
coerenti con contatori di sequenza o generazione. Un lettore lento può rilevare
di avere saltato aggiornamenti, ma non mantiene il produttore in attesa. Questa
regola è particolarmente rilevante quando telemetria e controllo vengono
eseguiti su contesti o core differenti.

La telemetria è suddivisa per significato:

\begin{itemize}
    \item il percorso di \emph{stato} contiene l'ultima osservazione disponibile;
    \item il percorso di \emph{eventi} conserva transizioni e richieste discrete;
    \item il percorso per \emph{rivoluzione} consente di ricostruire il moto e
          le decisioni ciclo per ciclo;
    \item il percorso \emph{diagnostico} contiene dati opzionali di sviluppo.
\end{itemize}

Quando la rete è congestionata, gli aggiornamenti ripetuti dello stato possono
essere ridotti. Eventi di sicurezza e lacune nei dati registrati devono invece
essere segnalati. Nessuna completezza del log giustifica il ritardo di un evento
motore.

\subsection{Controllo e sicurezza}

Il sottosistema di controllo riceve uno snapshot degli ingressi e uno stato di
sincronizzazione.
Una macchina a stati distingue almeno inizializzazione, acquisizione del
riferimento, funzionamento, taglio richiesto e condizione di allarme. Solo uno
stato con riferimento angolare attendibile può produrre una schedulazione CDI.

La mappa fornisce un valore base; correzioni termiche, gestione della
detonazione, limitatori e quick shifter producono richieste motivate. La
strategia le combina e registra sia il risultato richiesto sia quello
effettivamente comandato. Questa distinzione serve al digital twin: conoscere
un anticipo finale senza sapere se sia stato limitato non consente di spiegare
la decisione.

\subsection{Digital twin come estensione del sistema}

Il digital twin è collocato sul server e non sulla ECU. Per ogni dispositivo
mantiene identità hardware, cronologia del firmware, configurazioni, mappe,
calibrazioni, sessioni ed eventi. I dati vengono distinti in misurati, derivati,
decisi, comandati, stimati e simulati.

Il twin ricostruisce le sessioni senza alterare i campioni originali, riesegue
la logica di controllo sugli ingressi registrati e confronta l'uscita con mappe
o parametri alternativi. Ogni risultato mantiene la provenienza e la revisione
che lo ha prodotto. In questo modo il digital twin non coincide con un semplice
grafico di telemetria, ma costituisce uno strumento di replay, confronto e
diagnosi.

\section{Il prototipo}
\label{sec:prototipo}

\subsection{Sottosistemi realizzati}

Il prototipo integra tre catene funzionali. La prima acquisisce gli ingressi,
stima lo stato del motore, calcola anticipo e riferimenti ausiliari e temporizza
CDI, Power Jet e valvola di scarico. La seconda pubblica stato, eventi,
decisioni e comandi verso l'interfaccia web. La terza registra le sessioni e le
trasferisce al server, dove alimentano storico e digital twin.

\begin{center}
sensori $\rightarrow$ stato motore $\rightarrow$ controllo $\rightarrow$
attuatori

stato e decisioni $\rightarrow$ telemetria $\rightarrow$ client
$\rightarrow$ server
\end{center}

Le catene condividono tipi e identificativi, ma non le scadenze. Il percorso di
controllo termina sulla ECU; telemetria, registrazione e analisi consumano copie
coerenti e non possono ritardare la schedulazione degli output.

\subsection{Piattaforma di prototipazione}

La piattaforma è una scheda ESP32-S3 Mini basata su ESP32-S3FH4R2, con due core
Xtensa LX7, 4 MB di flash, 2 MB di PSRAM, Wi-Fi a 2,4 GHz e logica a 3,3 V. La
scheda rende disponibili ADC, GPIO, SPI, cattura temporizzata e connettività in
un formato adatto al banco. I dati elettrici e le periferiche derivano dalla
documentazione del produttore~\cite{esp32s3_datasheet}.

Il firmware è scritto in C++ sull'ambiente ESP-IDF e usa IDF FreeRTOS. La
documentazione Espressif descrive l'implementazione SMP dual-core e la possibilità
di assegnare affinità a uno specifico core~\cite{espidf_freertos}. Nel prototipo
questa capacità permette di tenere acquisizione e pubblicazione sensoriale
separate dal servizio HTTP/WebSocket. Priorità, sezioni critiche e carico sono
verificati nelle prove temporali, mentre le periferiche hardware eseguono
cattura e generazione degli eventi con intervento minimo della CPU.

Il banco usa circuiti di condizionamento per adattare i segnali ai livelli a
3,3 V della scheda. Alimentazione, protezioni e interfacce verso sensori e
attuatori restano esterne al modulo di calcolo, così da mantenere separata la
valutazione del firmware da quella elettrica.

\subsection{Ingressi e harness}

Ogni famiglia di ingressi è esposta al dominio attraverso una porta. La sorgente
può essere reale oppure sintetica; i servizi a valle non cambiano. Il profilo
di test usa sorgenti sintetiche deterministiche per TPS, temperature, pick-up,
quick shifter, selettore mappa e knock. Lo stesso percorso di validazione,
pubblicazione e controllo viene esercitato con ingressi reali o simulati.

Gli ingressi fisici sono:

\begin{itemize}
    \item TPS su un canale ADC, con conversione calibrata in millivolt;
    \item quick shifter su un ingresso digitale attivo basso;
    \item selettore di mappa su un ingresso digitale con debounce e validazione
          dello stato;
    \item pick-up già condizionato a onda quadra, acquisito mediante una
          periferica di cattura sul fronte di discesa;
    \item EGT tramite convertitore per termocoppia su bus SPI;
    \item temperatura acqua tramite partitore NTC e canale ADC;
    \item knock tramite interfaccia dedicata e finestra sincronizzata con la
          rivoluzione del motore.
\end{itemize}

\subsection{Modello sensoriale}

Una lettura contiene valore, timestamp monotono, validità, qualità, salute,
sequenza e bit di guasto applicabili. Il deposito sensoriale produce uno
snapshot coerente degli ingressi e mantiene code limitate per gli eventi. Il
TPS applica calibrazione e controlli di intervallo; il pick-up rifiuta intervalli
impossibili e perde la sincronizzazione quando gli eventi diventano obsoleti;
gli ingressi digitali implementano debounce e riarmo.

In presenza di un TPS non valido il servizio pubblica un fallback esplicito al
70 per cento. La strategia di sicurezza riconosce l'origine sostitutiva del
valore, applica i limiti previsti dal profilo operativo e registra la condizione
in telemetria.

Le temperature pubblicano stato e richieste di protezione, ma non eseguono lo
spegnimento. Analogamente, il sottosistema knock pubblica misure finestrate e
qualità; la strategia di accensione decide se e come ritardare
l'anticipo. Questa separazione conserva una sola autorità sul comando finale.

\subsection{Controllo motore e attuatori}

Gli eventi di pick-up aggiornano regime, sincronizzazione e identificativo di
rivoluzione. La macchina a stati abilita il controllo soltanto quando il
riferimento è attendibile. La mappa attiva fornisce l'anticipo base; correzioni
termiche, knock, limitatore e quick shifter producono il comando finale e un
codice che ne descrive la motivazione.

L'angolo viene convertito in un intervallo temporale rispetto al pick-up e
programmato su una periferica hardware. Power Jet e valvola di scarico usano
uscite PWM calcolate dalle rispettive mappe. Il sottosistema di attuazione
pubblica separatamente riferimento richiesto, comando applicato, saturazioni e
inibizioni, così che telemetria e digital twin possano ricostruire la decisione.

\subsection{Telemetria e interfaccia web}

Il componente di telemetria raccoglie snapshot ed eventi senza leggere
direttamente i driver. Il contratto espone stato dei sensori, sincronizzazione,
salute, contatori, riepilogo knock, eventi ordinati, decisioni di controllo e
comandi degli attuatori. Le informazioni destinate alla visualizzazione vengono
raggruppate in piccoli batch e serializzate in JSON.

JSON è più oneroso di un formato binario, ma rende il contratto leggibile e
semplifica l'ispezione con strumenti comuni. Per la dashboard il compromesso è
adeguato. I record per rivoluzione e le sessioni vengono invece codificati in
forma binaria compatta prima del trasferimento al server.

Il server HTTP/WebSocket gira sulla ECU. Pubblica un messaggio di capacità,
invia la telemetria quando il trasporto è pronto e non preleva nuovi batch
durante disconnessione o backpressure. Lo stesso server distribuisce gli asset
statici compressi dell'interfaccia. La Web UI adatta il contratto ai propri
modelli di visualizzazione, conserva uno storico locale limitato e mostra
letture, salute ed eventi recenti. I componenti dell'interfaccia non interpretano
direttamente il JSON ricevuto: il parser costituisce il confine tra protocollo e
presentazione. Dalla stessa interfaccia l'utente seleziona e modifica le mappe,
consulta le sessioni, avvia i comandi diagnostici e gestisce gli aggiornamenti
firmware. Le richieste vengono validate dalla ECU prima dell'applicazione.

\subsection{Mappe, sessioni e aggiornamenti}

Le mappe di accensione e Power Jet sono memorizzate in memoria non volatile e
usano punti RPM con interpolazione lineare. Il sistema verifica ordinamento,
intervalli e completezza prima di salvare una revisione. L'attivazione sostituisce
l'intera configurazione in un confine sicuro, evitando che il controllo osservi
una mappa modificata solo in parte.

La registrazione distingue la telemetria live dai dati di sessione. La prima è
ottimizzata per l'aggiornamento della dashboard; i secondi conservano record per
rivoluzione, eventi, decisioni e comandi in forma binaria. Il client inoltra i
blocchi al server insieme a identità ECU, firmware, mappa, configurazione e
calibrazione. Il server verifica completezza e sequenza prima di aggiornare il
digital twin.

L'aggiornamento OTA confronta la versione installata con l'artefatto approvato,
verifica integrità e autenticità dell'immagine e scrive la partizione inattiva.
Il nuovo firmware viene confermato dopo l'avvio corretto; in caso contrario il
bootloader ripristina la versione precedente.

\subsection{Scelte implementative rilevanti}

Quattro scelte caratterizzano l'implementazione.

La prima è l'uso di tempo iniettabile e sorgenti deterministiche. Timeout,
obsolescenza e debounce possono essere provati senza attendere il tempo reale.
La seconda è il limite esplicito di code e storico: un sovraccarico produce un
contatore o una perdita riconoscibile. La terza è l'assenza di dipendenze da
rete e storage nel dominio sensoriale. La quarta è la dichiarazione delle
capacità del protocollo, che permette a client e firmware di evolvere senza
assumere che ogni versione esponga gli stessi dati.

Le sorgenti vengono selezionate mediante profili validati, così che una
configurazione non possa associare un driver a un ingresso elettrico
incompatibile. JSON è riservato alla telemetria live, mentre la registrazione
usa un formato compatto e versionato. Le operazioni che modificano mappe o
firmware richiedono autenticazione e autorizzazione; gli aggiornamenti sono
verificati prima dell'avvio e prevedono il ripristino della versione precedente.

\section{Validazione}
\label{sec:validazione}

\subsection{Obiettivo della validazione}

La validazione deve rispondere a una domanda più stretta di ``il software
funziona'': quali requisiti sono sostenuti da un'evidenza ripetibile e a quale
livello? Sono considerati quattro livelli. I test sul calcolatore esercitano il
dominio senza hardware; la configurazione interamente sintetica esegue i servizi
sul dispositivo con ingressi deterministici; la modalità mista verifica i
singoli collegamenti fisici; il banco HIL misura il comportamento integrato di
ingressi, controllo, attuatori e telemetria.

\subsection{Test su calcolatore}

I test del dominio sensoriale eseguiti sul calcolatore coprono i casi seguenti:

\begin{itemize}
    \item conversione dei timestamp di cattura e gestione del ritorno del
          contatore;
    \item calibrazione TPS, intervalli invalidi, obsolescenza, fallback e
          recupero;
    \item debounce, riarmo e condizioni iniziali di quick shifter e map switch;
    \item stima RPM, acquisizione e perdita della sincronizzazione, intervalli
          impossibili e rapido aumento di regime;
    \item guasti termici, richieste di protezione e recupero senza comando
          diretto degli attuatori;
    \item misure knock mancanti, sature o arretrate;
    \item coerenza degli snapshot, sequenze e comportamento delle code piene;
    \item instradamento delle sorgenti sintetiche e reali attraverso l'harness;
    \item interpolazione delle mappe, composizione delle correzioni e
          transizioni della macchina a stati;
    \item validazione, persistenza e attivazione atomica delle mappe;
    \item inibizione dell'accensione con riferimento angolare non valido.
\end{itemize}

I test della telemetria verificano la composizione dello stato più recente,
l'ordinamento degli eventi e l'esposizione dei contatori di overflow. I test del
server verificano serializzazione, dichiarazione delle capacità, comportamento
in disconnessione e backpressure, risoluzione degli asset statici e rifiuto dei
percorsi non validi. Sul client vengono verificati parser, adattamento al modello
di visualizzazione, storico limitato, conservazione dell'ordine degli eventi e
validazione dei comandi. I test dei servizi remoti coprono registrazione delle
sessioni, associazione delle revisioni, immutabilità dei campioni, replay e
aggiornamento OTA con ripristino. L'esecuzione automatica associa i risultati a
revisione del codice, compilatore e configurazione usati nella prova.

\subsection{Prove sul dispositivo}

In modalità interamente sintetica il firmware esegue sulla ESP32-S3 e pubblica
periodicamente uno snapshot in formato CSV o attraverso la telemetria. La sequenza attesa è
specifica: il regime diventa valido dopo un numero sufficiente di catture, TPS e
temperature cambiano secondo gli stimoli, e gli eventi di quick shifter e cambio
mappa compaiono in punti ripetibili. Questa prova collega servizi, deposito e
runtime senza introdurre l'incertezza del cablaggio.

La modalità mista serve a verificare un confine elettrico alla volta. Una rampa
tra 0 e 3,3 V sull'ingresso analogico deve coprire l'intervallo calibrato del
TPS; l'ingresso del quick shifter deve produrre una sola richiesta validata
dopo debounce e riarmo; il selettore deve pubblicare il cambio della richiesta
fisica; un'onda quadra condizionata sull'ingresso di cattura produce il
regime coerente con la frequenza e diventa obsoleta quando gli impulsi
cessano. Le prove termiche verificano lettura nominale, termocoppia aperta,
guasto del convertitore e superamento delle soglie; la catena knock viene
esercitata con finestre valide, assenti e sature.

Il banco HIL applica profili di accelerazione, impulsi mancanti, disturbi, corto
verso massa o alimentazione e guasti dei convertitori. La strumentazione esterna
misura la differenza tra il riferimento angolare simulato e l'uscita CDI e
verifica duty cycle e frequenza degli attuatori ausiliari. Le stesse prove sono
ripetute con telemetria e trasferimento delle sessioni attivi, così da misurare
l'isolamento del percorso real-time.

\subsection{Tracciabilità rispetto ai requisiti}

La Tabella~\ref{tab:copertura} collega ogni requisito funzionale all'evidenza di
validazione corrispondente. I test sul calcolatore verificano la logica; le prove
sul dispositivo e HIL verificano i confini elettrici e temporali; le prove
end-to-end coprono client e server.

{\small
\begin{longtable}{p{0.14\textwidth} p{0.54\textwidth} p{0.18\textwidth}}
\caption{Copertura dei requisiti funzionali.}
\label{tab:copertura} \\
\hline
\textbf{Requisito} & \textbf{Evidenza di validazione} & \textbf{Stato} \\
\hline
\endfirsthead
\hline
\textbf{Requisito} & \textbf{Evidenza di validazione} & \textbf{Stato} \\
\hline
\endhead
\hline
\endfoot
\hline
\endlastfoot
RF-001 & Cattura del pick-up, stima RPM e acquisizione della sincronizzazione & Soddisfatto \\
RF-002 & Stima RPM, perdita e recupero della sincronizzazione, rapido aumento di regime & Soddisfatto \\
RF-003 & Controlli di intervallo, intervalli impossibili del pick-up e rifiuto dei valori non plausibili & Soddisfatto \\
RF-004 & Abilitazione del controllo solo con riferimento attendibile e inibizione con riferimento angolare non valido & Soddisfatto \\
RF-005 & Interpolazione della mappa, composizione delle correzioni e transizioni della macchina a stati & Soddisfatto \\
RF-006 & Profili validati delle sorgenti, validazione e persistenza dei parametri & Soddisfatto \\
RF-007 & Test dei sensori con sorgenti sintetiche e collegamenti fisici (modalità mista) & Soddisfatto \\
RF-008 & Calibrazione, controlli di intervallo e coerenza degli snapshot & Soddisfatto \\
RF-009 & Calibrazione TPS e validazione dei comandi di configurazione & Soddisfatto \\
RF-010 & Fallback TPS esplicito, obsolescenza e recupero & Soddisfatto \\
RF-011 & Persistenza e attivazione delle mappe, ripristino delle revisioni all'avvio & Soddisfatto \\
RF-012 & Validazione, persistenza e attivazione atomica delle mappe & Soddisfatto \\
RF-013 & Valori di sicurezza predefiniti, fallback espliciti e code piene & Soddisfatto \\
RF-014 & Validazione, applicazione e risposta ai comandi di configurazione & Soddisfatto \\
RF-015 & Viste, stato di salute, eventi e storico locale & Soddisfatto \\
RF-016 & WebSocket, sequenze, congestione e rilevamento delle lacune & Soddisfatto \\
RF-017 & Stato di salute dei sensori e contatori diagnostici & Soddisfatto \\
RF-018 & Obsolescenza, guasti, eventi di transizione, fallback e recupero & Soddisfatto \\
RF-019 & Prove termiche: lettura nominale, termocoppia aperta, guasto del convertitore e soglie & Soddisfatto \\
RF-020 & Pubblicazione dello stato termico e delle richieste di protezione & Soddisfatto \\
RF-021 & Misure knock mancanti, sature o arretrate e finestra sincronizzata & Soddisfatto \\
RF-022 & Quick shifter e selettore di mappa: debounce, riarmo e richieste validate & Soddisfatto \\
RF-023 & Cambio mappa e selezione della strategia in punti ripetibili & Soddisfatto \\
RF-024 & Persistenza, validazione e attivazione atomica delle mappe & Soddisfatto \\
RF-025 & Decisioni di controllo e codici di motivazione in telemetria & Soddisfatto \\
RF-026 & Comandi diagnostici, contatori e aggiornamento OTA con ripristino & Soddisfatto \\
RF-027 & Sorgenti sintetiche deterministiche e banco HIL con fault injection & Soddisfatto \\
RF-028 & Registrazione delle sessioni, immutabilità dei campioni e replay & Soddisfatto \\
RF-029 & Coerenza degli snapshot, telemetria e ricostruzione delle sessioni & Soddisfatto \\
\end{longtable}
}

\subsection{Valutazione dei requisiti non funzionali}

RNF-013 e RNF-014 sono verificati compilando il dominio senza ESP-IDF e
controllando il tempo mediante sorgenti deterministiche, così da rendere le
logiche di controllo testabili indipendentemente dall'hardware e le prove
ripetibili. RNF-019 è verificato saturando code e storici: in condizioni di
risorse limitate il sistema preserva le informazioni critiche ed espone contatori
e lacune senza crescita non limitata della memoria. RNF-010 è verificato
attraverso le regole di singolo scrittore e le istantanee coerenti, che mantengono
il comportamento del controllo indipendente dal carico delle altre funzioni. La
disconnessione di client e server conferma RNF-002 e RNF-006, poiché l'isolamento
delle funzioni critiche e la telemetria non bloccante lasciano invariata la
strategia di controllo attiva.

Per RNF-001, RNF-003 e RNF-004 le misure HIL confrontano latenza e jitter con il
budget temporale del controllo, sia a rete inattiva sia durante il massimo carico
di telemetria. Le prove di concorrenza confermano RNF-005, poiché acquisizione,
schedulazione e pubblicazione sensoriale non attendono i consumatori. La coerenza
degli snapshot verifica inoltre che ogni lettura conservi validità e salute
(RNF-016). RNF-018 e RNF-020 sono verificati ricostruendo una sessione dalla
relativa identità ECU, versione firmware, mappa, configurazione e calibrazione,
mantenendo ordine e provenienza dei dati. Le prove HIL con disturbi, corto verso
massa o alimentazione e guasti dei convertitori sostengono RNF-017, mentre il
fallback esplicito degli ingressi conferma la degradazione controllata richiesta
da RNF-023. L'aggiornamento OTA con verifica dell'artefatto, autorizzazione e
ripristino della versione precedente sostiene RNF-007, RNF-009 e RNF-022.

I risultati mostrano che le funzioni del sistema sono coperte da prove coerenti con
il loro livello: logica sul calcolatore, integrazione sul dispositivo, tempi e
segnali sul banco HIL, persistenza nelle prove end-to-end.

\section{Conclusioni e sviluppi futuri}
\label{sec:conclusioni}

\subsection{Risultati}

Il progetto definisce una ECU per motori monocilindrici come parte di un sistema
più ampio: sensori e attuatori sul veicolo, client locale, servizio remoto e
banco di simulazione. La separazione tra questi nodi stabilisce che il controllo
rimane sulla ECU, mentre interfaccia e server osservano, configurano e conservano
i dati senza entrare nel percorso real-time.

Il prototipo realizza l'intera catena descritta. Sorgenti reali o simulate
attraversano gli stessi servizi e producono snapshot ed eventi qualificati; il
controllo calcola e comanda accensione e attuatori; telemetria e Web UI rendono
osservabili misure e decisioni; il server conserva sessioni e revisioni e
fornisce replay e confronto mediante il digital twin. La validazione collega
queste funzioni ai requisiti e verifica l'isolamento del percorso real-time.

\subsection{Valutazione critica delle scelte}

La scheda ESP32-S3 offre un buon rapporto tra capacità di calcolo, periferiche e
connettività. Il front-end esterno mantiene condizionamento e protezioni fuori
dal modulo di calcolo; una scheda dedicata può in seguito integrare questi
elementi e ridurre cablaggio e ingombro. Il Wi-Fi rimane estraneo alle funzioni
essenziali e può essere disattivato senza interrompere il controllo.

Il fallback TPS fisso rende esplicita la condizione degradata e mantiene
deterministico il comportamento. Una politica dipendente dal punto operativo
può offrire limiti più conservativi, ma aumenta il numero di casi da calibrare
e verificare. Per questo il valore sostitutivo resta configurabile e viene
sempre accompagnato dallo stato di qualità.

Il contratto JSON mantiene leggibile la telemetria live, mentre il formato
binario versionato contiene il costo delle registrazioni per rivoluzione. La
doppia rappresentazione richiede un dizionario semantico comune, ma evita di
forzare lo stesso formato su visualizzazione e archiviazione. Il client come
gateway protegge l'autonomia della ECU; l'accodamento locale gestisce i periodi
nei quali il gateway non è disponibile.

Il digital twin mantiene separate osservazione, replay e simulazione. Questa
distinzione impedisce che un risultato ricalcolato venga interpretato come
misura fisica e rende confrontabili configurazioni diverse. I modelli
predittivi possono usare la stessa base dati, ma conservano una validazione
distinta dal controllo eseguito sulla ECU.

\subsection{Lavori futuri}

Gli sviluppi individuati estendono prestazioni, integrazione e capacità di
analisi senza modificare i confini principali del sistema:

\begin{itemize}
    \item mappe a due ingressi RPM--TPS con strumenti di confronto e
          ottimizzazione automatica delle superfici;
    \item strategie knock adattive che distinguano fondo meccanico, transitori
          e detonazione ripetuta;
    \item ampliamento del banco HIL con profili derivati da sessioni reali e
          campagne automatiche di fault injection;
    \item analisi temporale estesa a diverse condizioni di carico e temperatura;
    \item scheda dedicata con protezioni, condizionamento, alimentazione e
          connettori adatti all'ambiente del veicolo;
    \item integrazione con reti CAN e strumenti diagnostici automotive;
    \item modelli termici e predittivi aggiuntivi nel digital twin;
    \item campagne di calibrazione su motori con cilindrata e configurazione
          differenti.
\end{itemize}

Queste estensioni riusano gli stessi contratti tra acquisizione, controllo,
telemetria e analisi, mantenendo verificabili le nuove funzioni in modo
indipendente.

\begin{thebibliography}{9}

\bibitem{linkecu}
Link Engine Management,
\emph{Why an Aftermarket ECU? Selecting an Aftermarket Engine Management
System}.
\url{https://linkecu.com/new-to-ecus/why-an-aftermarket-ecu/}.

\bibitem{iea_sdv}
International Energy Agency,
\emph{Vehicle software and software-defined vehicles},
20 maggio 2026.
\url{https://www.iea.org/reports/vehicle-software-and-software-defined-vehicles}.

\bibitem{autosar_classic}
AUTOSAR,
\emph{Classic Platform}.
\url{https://www.autosar.org/standards/classic-platform}.

\bibitem{abboush_hil}
M. Abboush, C. Knieke e A. Rausch,
\emph{A Virtual Testing Framework for Real-Time Validation of Automotive
Software Systems Based on Hardware in the Loop and Fault Injection},
Sensors, vol. 24, n. 12, 2024, doi: 10.3390/s24123733.

\bibitem{tran_dt}
V. D. Tran, P. Sharma e L. H. Nguyen,
\emph{Digital Twins for Internal Combustion Engines: A Brief Review},
Journal of Emerging Science and Engineering, vol. 1, n. 1, pp. 29--35,
2023, doi: 10.61435/jese.2023.5.

\bibitem{nath_dt}
K. Nath, V. Kumar, D. J. Smith e G. E. Karniadakis,
\emph{A Digital Twin for Diesel Engines: Operator-infused Physics-Informed
Neural Networks with Transfer Learning for Engine Health Monitoring},
Engineering Applications of Artificial Intelligence, 2026,
doi: 10.1016/j.engappai.2026.114052.

\bibitem{esp32s3_datasheet}
Espressif Systems,
\emph{ESP32-S3 Series Datasheet}.
\url{https://documentation.espressif.com/esp32-s3_datasheet_en.pdf}.

\bibitem{espidf_freertos}
Espressif Systems,
\emph{FreeRTOS (IDF) -- ESP-IDF Programming Guide}.
\url{https://docs.espressif.com/projects/esp-idf/en/stable/esp32s3/api-reference/system/freertos_idf.html}.

\end{thebibliography}

\end{document}
